\chapter{اللقاء الثاني والعشرون: تتمة آيات القصاص والوصية وبداية آيات الصيام}

\section{مقدمة: مراجعة منهج الطبري في آية القصاص}
السلام عليكم ورحمة الله وبركاته. الحمد لله رب العالمين، وأشهد أن لا إله إلا الله وحده لا شريك له، وأشهد أن محمداً عبده ورسوله. اللهم صل على محمد وعلى آل محمد كما صليت على آل إبراهيم إنك حميد مجيد، اللهم بارك على محمد وعلى آل محمد كما باركت على آل إبراهيم إنك حميد مجيد.

افتتح الشيخ هذا اللقاء بمراجعة الواجب السابق في تلخيص منهج الإمام \rawi{الطبري} في آية القصاص، وبيّن من خلال مناقشة الطلاب كيف يجمع الطبري الأقوال، ويميز بين الخلاف في سبب النزول والخلاف في التفسير، ثم يرد المسألة إلى الحكم المحكم بدلالة النصوص والإجماع.

\begin{fawaid}[الفرق بين الخلاف في النزول والخلاف في التفسير]
ليس كل اختلاف يذكره المفسرون في الآية اختلافاً في معناها المراد ابتداءً، بل قد يكون اختلافاً في سبب نزولها، أو في بعض أفرادها، أو في تطبيقها. والتمييز بين هذه الأنواع يعين على فهم طريقة الطبري في الترجيح وعدم الخلط بين المراتب.
\end{fawaid}

\separator

\section{تأويل (ذلك تخفيف من ربكم ورحمة)}
\isnad{قال أبو جعفر:} هذا الذي شرعه الله للأمة في أمر القصاص، من إباحة العفو إلى الدية، تخفيف من الله ورحمة بهذه الأمة، بخلاف من قبلهم ممن كان الحكم فيهم أضيق.

وذكّر الشيخ هنا بما مضى من أن من أعظم معاني الشريعة رحمة الله بعباده، وأن التشريع ليس مجرد إلزام ومشقة، بل هو عدل ورحمة وتيسير وحفظ لمصالح الخلق.

\begin{fawaid}[الشريعة رحمة لا مجرد تكليف مجرد]
من الخطأ اختزال أحكام الشريعة في صورة الأوامر الشاقة المجردة؛ بل من أعظم مقاصدها الرحمة والتخفيف وحفظ الدين والنفس والمال والعرض والعقل. وهذا المعنى يظهر كثيراً في تعقيبات القرآن على الأحكام.
\end{fawaid}

\separator

\section{تأويل (فمن اعتدى بعد ذلك فله عذاب أليم)}
ذكر \rawi{الطبري} الأقوال فيمن المراد بهذا المعتدي، ثم رجح أن من أخذ الدية أو رضي بالعفو ثم عاد بعد ذلك إلى قتل القاتل أو التعدي، فقد اعتدى وظلم، وله عذاب أليم.

ونبّه الشيخ في أثناء الشرح إلى أن هذا الترجيح يكشف عن المذهب الفقهي الذي يفهمه \rawi{الطبري} من الآية: وهو أن سلطان ولي الدم باقٍ في حدود ما شرعه الله، ولا يجوز له بعد الصلح أو أخذ الدية أن يتجاوز ذلك.

\separator

\section{تأويل (ولكم في القصاص حياة يا أولي الألباب)}
\isnad{قال أبو جعفر:} في القصاص حياة للناس؛ لأن من علم أنه متى قتل قُتل، انكف عن القتل، فكان في ذلك بقاء للنفوس وردع للعدوان.

وساق عن \rawi{مجاهد} و\rawi{قتادة} و\rawi{الربيع} وغيرهم أن معنى الحياة هنا: النكال والمنع والزجر، وأن الله حجز بعض الناس عن بعض بالقصاص.

\begin{fawaid}[الزجر والعقوبة من مقاصد الشريعة في حفظ المجتمع]
لا تكتفي الشريعة بمعاقبة الجاني بعد وقوع الجريمة، بل تجعل العقوبة نفسها زاجرةً رادعةً تمنع وقوع الجريمة ابتداءً. وهذا من معنى قوله تعالى: \textit{وَلَكُمْ فِي الْقِصَاصِ حَيَاةٌ}، وهو أصل في أن العقوبات الشرعية من وسائل حفظ الحياة لا مناقضتها.
\end{fawaid}

\begin{fawaid}[التقوى في كل آية تفسر بحسب سياقها أولاً]
قوله تعالى: \textit{لَعَلَّكُمْ تَتَّقُونَ} وإن كان يتناول التقوى عموماً، فإن أول ما يحمل عليه هو ما يقتضيه السياق. فهنا: لعلكم تتقون القتل المؤدي إلى القصاص، ثم يدخل فيه بعد ذلك ما يلحق به من المعاني العامة.
\end{fawaid}

\separator

\section{تأويل (كتب عليكم إذا حضر أحدكم الموت إن ترك خيراً الوصية)}
\isnad{قال أبو جعفر:} فُرض عليكم إذا حضر أحدكم الموت وترك مالاً أن يوصي للوالدين والأقربين الذين لا يرثونه بالمعروف، من غير مجاوزة للثلث ولا ظلم للورثة.

وأطال \rawi{الطبري} في مناقشة دعوى النسخ هنا، وذكر قول من قال إن الآية منسوخة بآيات المواريث وحديث: \textit{لا وصية لوارث}، وذكر في المقابل قول من قال إنها محكمة، وأن المراد بها من لا يرث من الوالدين والأقربين.

ثم رجح \rawi{الطبري} أن الآية محكمة غير منسوخة، وأنه لا يجوز ادعاء النسخ ما دام يمكن اجتماع حكم الوصية مع حكم المواريث على وجه صحيح؛ بأن تكون الوصية لمن لا يرث.

\begin{fawaid}[تضييق دائرة النسخ عند الطبري]
من أصول الطبري أنه لا يحكم بالنسخ لمجرد وجود قولين أو ظاهر تعارض، بل يشترط أن يكون الحكم الثاني نافياً للأول على وجه لا يمكن معه اجتماعهما في حال واحدة. فإذا أمكن الجمع الصحيح بين الآيتين لم يجز دعوى النسخ.
\end{fawaid}

\begin{fawaid}[لا يصرف الظاهر إلى باطن بلا حجة]
رجح الطبري وجوب الوصية في قليل المال وكثيره، لأن ظاهر الآية يقتضي ذلك، ولم يثبت عنده ما يخصص المال الكثير دون القليل. وهذه صورة أخرى من تقديم ظاهر التنزيل حتى تقوم حجة صحيحة على التخصيص.
\end{fawaid}

\subsection{تأويل (إن ترك خيراً)}
فسر \rawi{الطبري} الخير هنا بالمال، وذكر الخلاف في تحديد القليل والكثير منه، ثم رجح أن كل مال حضر الموت صاحبه وعنده منه شيء فهو داخل في الآية، قليلاً كان أو كثيراً، ما دام يصدق عليه اسم المال.

\subsection{تأويل (فمن بدله بعد ما سمعه)}
بيّن \rawi{الطبري} أن الإثم على من بدّل الوصية المشروعة بعد سماعها، لا على الموصي إذا كان قد أوصى بالحق، لأن الله سميع لوصيته عليم بنيات عباده.

\subsection{تأويل (فمن خاف من موص جنفاً أو إثماً)}
اختلف أهل التأويل في الجنف والإثم هنا، وذكر \rawi{الطبري} أن من حضر الموصي فرآه مائلاً عن الحق أو ظالماً في وصيته، فأصلح بينه وبين ورثته أو من أوصى لهم، فلا إثم عليه.

\begin{fawaid}[الإصلاح في الوصايا من إقامة العدل لا من تبديل الشرع]
إذا ظهر من الموصي ميل أو ظلم أو جنف، فإصلاح ذلك وردّه إلى العدل ليس تبديلًا للوصية المحقة، بل هو إصلاح مشروع أمر الله به، ومن قام به فلا إثم عليه.
\end{fawaid}

\separator

\section{تأويل (يا أيها الذين آمنوا كتب عليكم الصيام)}
\isnad{قال أبو جعفر:} فرض عليكم الصيام كما فرض على من قبلكم، رجاء أن تتقوا الله بامتثال أمره وترك معصيته.

وذكر الشيخ أن \rawi{الطبري} يبدأ دائماً من المعنى الذي نزلت فيه الآية أولاً، ثم ينظر في عمومها بعد ذلك، فلا يهمل السياق، ولا يحصر الدلالة فيه وحده.

\subsection{تأويل (أياماً معدودات)}
فسر \rawi{الطبري} المعدودات بأنها الأيام المحصيات المحدودات العدد.

\subsection{تأويل (فمن كان منكم مريضاً أو على سفر فعدة من أيام أخر)}
بيّن \rawi{الطبري} أن الله رخّص للمريض والمسافر في الفطر، وأوجب عليهما القضاء من أيام أخر، ثم دخل في تفصيل الخلاف في قوله تعالى: \textit{وعلى الذين يطيقونه فدية طعام مسكين}.

\separator

\section{تأويل (وعلى الذين يطيقونه فدية طعام مسكين)}
ذكر \rawi{الطبري} في هذه الآية أقوالاً كثيرة:
\begin{itemize}
    \item قول من قال: كانت رخصة عامة في أول الأمر، فمن شاء صام ومن شاء أفطر وافتدى، ثم نسخ ذلك بقوله: \textit{فَمَن شَهِدَ مِنكُمُ الشَّهْرَ فَلْيَصُمْهُ}.
    \item وقول من خصها بالشيخ الكبير والعجوز ومن في معناهما.
    \item وقول من أدخل فيها الحامل والمرضع.
    \item وقول من قرأها على معنى: \textit{يُطَوَّقونه} أي: يكلفونه مع المشقة.
\end{itemize}

ثم رجح \rawi{الطبري} أن الآية منسوخة في حق القادر الصحيح المقيم، وأنها كانت في أول الأمر تخييراً بين الصوم والفدية، ثم أوجب الله الصيام على من شهد الشهر.

\begin{fawaid}[تعليق صحة التفسير أحياناً على صحة الخبر]
من أمثلة دقة الطبري أنه قد يذكر خبراً يستدل به بعض أهل العلم على معنى في الآية، ثم يقول: ``إن صح الحديث'' كان معناه كذا. فهو لا يبني التفسير الجازم على خبر لم يثبت عنده، بل يعلّق الاستدلال على ثبوته.
\end{fawaid}

\subsection{حديث الحامل والمرضع}
ذكر \rawi{الطبري} الحديث الوارد في وضع الصوم عن الحامل والمرضع، ثم قال: إن صح فمعناه أنهما تفطران حال العجز ثم تقضيان، لا أن عليهما الفدية فقط من غير قضاء. وبذلك وافق أصول الآية وسياقها العام في القضاء.

\separator

\section{تأويل (ومن تطوع خيراً فهو خير له)}
ذكر \rawi{الطبري} الأقوال في معنى التطوع هنا، وهل المراد به الزيادة في الإطعام أو ما هو أعم من ذلك، وقرر أن كل ما زاده العبد من الخير فهو خير له عند الله.

\section{تأويل (وأن تصوموا خير لكم)}
فسرها \rawi{الطبري} بأن الصوم خير لكم من الفدية إذا قدرتم عليه، لأن الصوم هو الأصل الذي كتبه الله على عباده.

\separator

\section{تأويل (شهر رمضان الذي أنزل فيه القرآن)}
فسر \rawi{الطبري} الشهر، وتكلم على معنى إنزال القرآن في رمضان، وأنه هدى للناس وبينات من الهدى والفرقان، أي: فاصل بين الحق والباطل.

\section{تأويل (فمن شهد منكم الشهر فليصمه)}
ذكر \rawi{الطبري} الخلاف في معنى شهود الشهر: هل هو حضور المقيم، أو رؤية الهلال، أو بلوغ الشهر مع الأهلية للصوم؟ ثم رجح من المعاني ما يندرج تحت إلزام من أدرك الشهر وهو من أهل الصيام أن يصومه.

\separator

\section{تأويل (ومن كان مريضاً أو على سفر فعدة من أيام أخر)}
عاد \rawi{الطبري} هنا إلى تفصيل حكم الفطر في المرض والسفر، وذكر الخلاف: هل الإفطار عزيمة واجبة أو رخصة؟ ثم رجح أن الإفطار رخصة لا عزيمة، وأن للمريض والمسافر أن يفطرا، ويقضيا من أيام أخر.

\begin{fawaid}[الترجيح بالنظائر من القرآن واللغة]
من منهج الطبري في الترجيح أنه يستعمل النظائر من القرآن ومن أساليب العرب. وقد نبه الشيخ هنا إلى أن النظائر قد تكون في المعنى، وقد تكون في الأسلوب والتركيب، وكلاهما يفيد في ترجيح المراد من الآية.
\end{fawaid}

\separator

وقف الشيخ عند هذا الموضع من آيات الصيام، بعد أن بيّن جملة من مسائل النسخ والرخصة والعموم والخصوص في آيات الوصية والصيام، وربط ذلك كله بمنهج \rawi{الطبري} في الجمع بين الآثار، وتضييق النسخ، والرجوع إلى ظاهر التنزيل حتى تقوم الحجة بخلافه.
